% 字段与历史的三种融合位置；全版心，约162mm。
\begin{tikzpicture}[
 font=\figurefont\footnotesize,
 box/.style={draw=BookRule,fill=BookPaper,align=center,
   text width=32mm,minimum height=11mm,inner sep=3pt},
 flow/.style={-{Stealth[length=2mm]},draw=BookMuted,line width=.9pt}
]
 \node[font=\figurefont\footnotesize\bfseries,anchor=west] at (-1.8,1.35)
   {末端融合：先独立压缩，再比较候选};
 \node[box,fill=FigureInput!10] (h1) at (0,.55) {原始历史};
 \node[box,fill=FigureInput!10] (f1) at (0,-.85) {用户、候选与情境};
 \node[box,fill=FigureModel!10] (he1) at (4.8,.55) {历史表示};
 \node[box,fill=FigureModel!10] (fi1) at (4.8,-.85) {字段表示};
 \node[box,fill=FigureOutput!10] (o1) at (12,0) {融合及响应头};
 \draw[flow] (h1)--(he1);
 \draw[flow] (f1)--(fi1);
 \draw[flow] (he1.east)--(8,.55)--(o1.west);
 \draw[flow] (fi1.east)--(8,-.85)--(o1.west);
 \node[font=\figurefont\footnotesize\bfseries,anchor=west] at (-1.8,-2)
   {逐层桥接：两套表示通过摘要交换信息};
 \node[box,fill=FigureModel!10] (h2) at (0,-2.85) {历史模块};
 \node[box,fill=FigureModel!10] (f2) at (5.5,-2.85) {字段模块};
 \node[box,fill=FigureOutput!10] (o2) at (12,-2.85) {响应头};
 \draw[flow] (h2.east) to[bend left=20] node[above]{历史摘要} (f2.west);
 \draw[flow] (f2.west) to[bend left=20] node[below]{字段摘要} (h2.east);
 \draw[flow] (f2)--(o2);
 \node[font=\figurefont\footnotesize\bfseries,anchor=west] at (-1.8,-4.3)
   {联合骨干：每层同时安排字段交互与历史读取};
 \node[box,fill=FigureInput!10] (i3) at (0,-5.2) {历史与字段单元};
 \node[box,fill=FigureModel!10] (b3) at (5.5,-5.2) {联合层堆栈\\掩码控制信息可见};
 \node[box,fill=FigureOutput!10] (o3) at (12,-5.2) {候选响应头};
 \draw[flow] (i3)--(b3);
 \draw[flow] (b3)--(o3);
\end{tikzpicture}
